\documentclass[11pt]{amsart}
\usepackage[T1]{fontenc}
\usepackage{lmodern}
\usepackage[expansion=false]{microtype}
\usepackage[margin=1.1in]{geometry}
\usepackage{amsmath,amssymb,amsthm}
\usepackage[hidelinks]{hyperref}
\urlstyle{same}

\newcommand{\R}{\mathbb R}
\newcommand{\C}{\mathbb C}
\newcommand{\F}{\mathcal F}
\newcommand{\Core}{\mathcal C}
\newcommand{\Hphys}{\mathcal H}
\newcommand{\Spin}{\operatorname{Spin}}
\newcommand{\SU}{\operatorname{SU}}
\newcommand{\SO}{\operatorname{SO}}
\newcommand{\tr}{\operatorname{tr}}
\newcommand{\supp}{\operatorname{supp}}
\newcommand{\Dom}{\operatorname{Dom}}
\newcommand{\ran}{\operatorname{ran}}
\newcommand{\norm}[1]{\lVert #1\rVert}
\newcommand{\ip}[2]{\langle #1,#2\rangle}

\newtheorem{theorem}{Theorem}[section]
\newtheorem{proposition}[theorem]{Proposition}
\newtheorem{lemma}[theorem]{Lemma}
\newtheorem{corollary}[theorem]{Corollary}
\theoremstyle{remark}
\newtheorem{remark}[theorem]{Remark}
\numberwithin{equation}{section}

\title[Positive eigenvalues of the relative $\SU(2)$ BFSS Hamiltonian]
{Positive eigenvalues of the relative $\SU(2)$ BFSS Hamiltonian}
\author{OpenAI}
\date{October 5, 2026}
\hypersetup{
  pdftitle={Positive eigenvalues of the relative SU(2) BFSS Hamiltonian},
  pdfauthor={OpenAI}
}

\begin{document}

\begin{abstract}
We prove that the relative $\SU(2)$ BFSS Hamiltonian has infinitely many
positive eigenvalues tending to infinity, with square-integrable
eigenvectors. The operator is defined by closing the gauge-invariant
supercharge form. This refutes, at $N=2$, the exclusion of normalizable
positive-energy states stated in the original BFSS paper.
\end{abstract}

\maketitle

\section{Introduction}
\label{sec:introduction}

The Banks--Fischler--Shenker--Susskind (BFSS) model is supersymmetric
matrix quantum mechanics obtained by reducing ten-dimensional
supersymmetric Yang--Mills theory to one time dimension. Its quartic
potential vanishes on commuting matrices, leaving unbounded directions
of classical escape. For the relative $\SU(2)$ Hamiltonian defined
below, we prove an unbounded sequence of positive eigenvalues with
square-integrable eigenvectors. The proof selects $\Spin(9)$ sectors
in which symmetry excludes low transverse oscillator levels, so the
remaining oscillator energy confines those sectors. We first specify
the operator and its physical Hilbert space.

We study the relative model, with the free center of mass removed, and use
the following normalization. The color basis $T_a=\sigma_a/\sqrt2$,
where $\sigma_a$ are the Pauli matrices, satisfies
\[
 \operatorname{Tr}(T_aT_b)=\delta_{ab},\qquad
 [T_a,T_b]=i f_{abc}T_c,\qquad f_{abc}=\sqrt2\,\varepsilon_{abc}.
\]
Write the nine traceless bosonic matrices as
$X_i=\sum_{a=1}^3x_i^aT_a$ and group their coefficients into
\[
 x=(x^1,x^2,x^3)\in(\R^9)^3,\qquad x^a=(x_i^a)_{i=1}^9,\qquad
 |x|^2=\sum_{a=1}^3|x^a|^2.
\]
An orthogonal change of color basis acts unitarily on the bosonic and
fermionic variables, so this choice represents the $\SU(2)$ Hamiltonian
in the stated normalization.

Let $\F$ be an irreducible complex Clifford module with self-adjoint
generators $\theta_\alpha^a$, $1\leq\alpha\leq16$ and $1\leq a\leq3$,
such that
\[
 \{\theta_\alpha^a,\theta_\beta^b\}
   =\delta_{\alpha\beta}\delta_{ab}.
\]
In particular, $\F$ is finite dimensional. Fix real symmetric matrices
$\gamma^1,\ldots,\gamma^9$ of size $16$ with
\[
 \{\gamma^i,\gamma^j\}=2\delta_{ij}I,\qquad
 \gamma^{ij}=\tfrac12[\gamma^i,\gamma^j].
\]
For $h\in\SU(2)$, let $O_h\in\SO(3)$ be its adjoint action in the
basis $T_a$. Its action on the real generator label space
$\R^3\otimes\R^{16}$ is $O_h\otimes I_{16}\in\SO(48)$.
Since $\SU(2)$ is simply connected, this homomorphism lifts uniquely
to $\Spin(48)$. Let $V_h$ be the unitary Clifford action of that lift
on $\F$, using the normalized generators $\sqrt2\,\theta_\alpha^a$.
The gauge action on functions is
\[
 (\mathcal G_h\Psi)(x)=V_h\Psi(O_h^{-1}x),
\]
where $O_h^{-1}$ acts on the color index of $x$. Define
\[
 \Hphys=\bigl[L^2(\R^{27};\F)\bigr]^{\SU(2)},\qquad
 \Core=\bigl[C_c^\infty(\R^{27};\F)\bigr]^{\SU(2)}.
\]
On $\Core$, put
\begin{equation}
 Q_\alpha
 =-i\gamma^i_{\alpha\beta}\theta_\beta^a\partial_{ia}
   +\tfrac12 f_{abc}x_i^b x_j^c
      \gamma^{ij}_{\alpha\beta}\theta_\beta^a,
 \qquad
 q(\Psi)=\tfrac1{16}\sum_{\alpha=1}^{16}\norm{Q_\alpha\Psi}_2^2.
 \label{eq:supercharges}
\end{equation}
Here $\partial_{ia}=\partial/\partial x_i^a$, and repeated indices are
summed. The tensors $\delta_{ab}$ and $f_{abc}$ are invariant under
$O_h$, so $\mathcal G_hQ_\alpha\mathcal G_h^{-1}=Q_\alpha$ on
$C_c^\infty(\R^{27};\F)$. Each $Q_\alpha$ is symmetric on $\Core$.
Gauge averaging shows that $\Core$ is dense in $\Hphys$.
Consequently the operator
$\Psi\mapsto(Q_1\Psi,\ldots,Q_{16}\Psi)$ is closable: convergence of
$\Psi_n$ to zero and of each $Q_\alpha\Psi_n$ forces every limit to vanish
by symmetry. Thus $q$ is closable. Let $H$ be the nonnegative self-adjoint
operator associated with its closure, still denoted by $q$.
This is the relative $\SU(2)$ BFSS Hamiltonian considered here.

\begin{theorem}
\label{thm:main}
There are an orthonormal sequence
$\Psi_j\in\Dom(H)$ and real numbers $E_j>0$ tending to infinity such that
\[
 H\Psi_j=E_j\Psi_j.
\]
In particular, $\sigma_{\mathrm p}(H)\cap(0,\infty)$ is infinite.
\end{theorem}

The mechanism is visible in the quartic bosonic potential
\[
 V(x)=\sum_{a<b}|x^a\wedge x^b|^2.
\]
It vanishes when the three spatial vectors are collinear, leaving an
unbounded set along which confinement can fail. Fixing one nonzero vector
$u=x^a$ separates each other vector into a component parallel to $u$ and
a component in $u^\perp$. The potential together with the transverse
kinetic energy bounds a harmonic oscillator in these sixteen
coordinates, with frequency proportional to $|u|$.

The symmetry determines which oscillator levels are available. The
stabilizer of $u$ in $\Spin(9)$ is $\Spin(8)$. If a $\Spin(9)$
representation has highest weight
$\lambda=(\lambda_1,\lambda_2,\lambda_3,\lambda_4)$, every constituent
of its restriction to this stabilizer has highest weight $\nu$ with
$\nu_1\geq\lambda_2$. If $n$ is the oscillator level and $M$ is a
fixed bound for the contribution from the fermions, a stabilizer type
$\nu$ shared by this restriction and level $n$ tensored with the
fermions must satisfy
\[
 \lambda_2\leq\nu_1\leq n+M.
\]
A large $\lambda_2$ therefore removes low levels. This is why the second
coordinate matters: the first coordinate alone does not impose this
restriction on stabilizer types.

The resulting oscillator energy grows linearly in $|u|$. Squaring the
supercharges also produces a fermionic multiplication term that is linear
in $x$ and may be negative. By excluding sufficiently many oscillator
levels, the positive linear coefficient dominates this term. Averaging
the three choices of color then gives a form bound that confines the
full variable $x$ within the chosen rotation sector. For every sector
above a fixed threshold for $\lambda_2$, the restricted Hamiltonian
consequently has compact resolvent and trivial kernel. Explicit
gauge-invariant polynomials
show that infinitely many such physical sectors are infinite
dimensional; any one of them supplies the unbounded sequence of
positive eigenvalues.

The spectral history makes the distinction between the spectral set
and normalizable eigenstates essential. For the finite matrix
Hamiltonians they studied, de Wit, L\"uscher, and Nicolai proved that
the spectral set is $[0,\infty)$ and explicitly left open the
possibility of normalizable eigenstates within that interval
\cite[Section~1]{deWitLuscherNicolai}. In their discussion of relative
motion after removing the free center of mass, Banks, Fischler, Shenker,
and Susskind stated that, besides the zero-energy threshold states,
\textquotedblleft No other normalizable bound states can occur\textquotedblright~
\cite[Section~5, after Equation~(5.4)]{BFSS}. Theorem~\ref{thm:main}
refutes this exclusion for the exact relative operator defined here
at $N=2$; the large-$N$ Matrix Theory conjecture is a separate question.

Positive-energy binding had appeared in a qualified form in the work
of Danielsson, Ferretti, and Sundborg. They obtained normalizable
excited wavefunctions in a Born--Oppenheimer approximation to the same
relative $\SU(2)$ model, anticipated decay beyond the approximation,
and described the states as metastable
\cite[Sections~3.3 and 4--5]{DanielssonFerrettiSundborg}. Smilga later
conjectured families of normalized excited states in a large-$N$
thermodynamic analysis \cite[Section~3]{Smilga}. Sethi and Stern's
rotational invariance theorem concerns normalizable zero-energy
states \cite[Sections~3.2 and 4.2]{SethiSternInvariance}, whereas the
eigenvalues proved to exist here are positive.

Transverse quantization has classical scalar precedents in Simon's
$x^2y^2$ oscillator argument \cite[Section~2]{Simon} and L\"uscher's
estimates for bosonic Yang--Mills Hamiltonians
\cite[Appendix~D]{Luscher}. Davies and Simon used parity to remove the
transverse constant mode of a symmetric Neumann horn and obtain a reducing sector
with compact resolvent \cite[Theorem~3.1 and Step~7]{DaviesSimon}.
In a supersymmetric matrix model with two spatial matrices and
$\SU(2)\times\SO(2)$ symmetry, Aref'eva, Koshelev, and Medvedev
decomposed the problem into angular sectors and used a
Born--Oppenheimer analysis to argue for sectors with purely discrete
spectrum, deriving asymptotic spectral formulas
\cite[Sections~5--6]{ArefevaKoshelevMedvedev}.
Here the decisive model-specific step is to use orthogonal branching
and the finite fermionic weight set to exclude any prescribed finite
number of transverse oscillator levels. More generally, restriction
to the stabilizer of an escape direction can identify symmetry sectors
in which transverse excitation restores confinement.

\section{The quadratic form and rotations}
\label{sec:form}

We first isolate the scalar potential and the fermionic term. We then
show that rotation types reduce the closed form and its associated
operator. All differential computations in this section take place on
compactly supported smooth functions.

\begin{lemma}
\label{lem:form}
For $\Psi\in\Core$,
\begin{equation}
 q(\Psi)=\tfrac12\norm{\nabla\Psi}_2^2
   +\int_{\R^{27}}V(x)\norm{\Psi(x)}^2\,dx
   +\int_{\R^{27}}\ip{\Psi(x)}{B(x)\Psi(x)}\,dx,
 \label{eq:form}
\end{equation}
where
\begin{equation}
 V(x)=\sum_{a<b}|x^a\wedge x^b|^2,\qquad
 B(x)=\tfrac i2 f_{abc}x_i^a\gamma^i_{\beta\delta}
                   \theta_\beta^b\theta_\delta^c.
 \label{eq:potential}
\end{equation}
The matrix $B(x)$ is Hermitian. There is a constant $C<\infty$ such that
\begin{equation}
 \norm{B(x)}_{\mathrm{op}}\leq C|x|\qquad(x\in\R^{27}).
 \label{eq:fermion-bound}
\end{equation}
\end{lemma}

\begin{proof}
The calculation also holds on $C_c^\infty(\R^{27};\F)$ before taking
gauge invariants. Put
\[
 y^a_{ij}=f_{abc}x_i^b x_j^c,\qquad
 A_\alpha=-i\gamma^i_{\alpha\beta}\theta_\beta^a\partial_{ia},
 \qquad
 W_\alpha=\sum_{j<k,c}y^c_{jk}\gamma^{jk}_{\alpha\delta}\theta_\delta^c.
\]
The antisymmetry in $j,k$ makes $Q_\alpha=A_\alpha+W_\alpha$.
The gamma relations give
\[
 \tr((\gamma^i)^T\gamma^j)=16\delta_{ij},\qquad
 \tr((\gamma^{ij})^T\gamma^{kl})
   =16\delta_{ik}\delta_{jl}\quad(i<j,\ k<l).
\]
Indeed, each indicated matrix has squared Hilbert--Schmidt norm $16$,
and a product of two or four distinct gamma matrices has trace zero:
conjugation by one of its factors changes its sign. Clifford
anticommutation, followed by these identities, gives on the smooth core
\[
 \tfrac1{16}\sum_\alpha A_\alpha^2=-\tfrac12\Delta,\qquad
 \tfrac1{16}\sum_\alpha W_\alpha^2
    =\tfrac12\sum_{a,j<k}(y^a_{jk})^2I.
\]
For the latter identity, the coefficients $y^a_{jk}$ commute, so the
two orders of each Clifford product combine into an anticommutator.
Since $f_{abc}=\sqrt2\,\varepsilon_{abc}$, the last scalar is
$\sum_{a<b}|x^a\wedge x^b|^2$.

In the cross term, the first-order part is
\[
 -\tfrac i{16}\sum_{i,a,j<k}
 y^a_{jk}\tr((\gamma^i)^T\gamma^{jk})\partial_{ia}=0,
\]
because $\gamma^i$ is symmetric and $\gamma^{jk}$ is skew-symmetric.
The remaining multiplication operator is
\[
 -\tfrac i{16}\sum_{\substack{i,a,\beta,\delta\\j<k,c}}
   (\partial_{ia}y^c_{jk})
   (\gamma^i\gamma^{jk})_{\beta\delta}
   \theta_\beta^a\theta_\delta^c.
\]
Here
\[
 \partial_{ia}y^c_{jk}
 =f_{cab}\delta_{ij}x_k^b+f_{cba}x_j^b\delta_{ik},
 \qquad
 \gamma^j\gamma^{jk}=\gamma^k,\quad
 \gamma^k\gamma^{jk}=-\gamma^j.
\]
Each spatial index occurs in eight pairs $j<k$. Substitution gives
\[
 -\tfrac i2 f_{cab}x_i^b\gamma^i_{\beta\delta}
                   \theta_\beta^a\theta_\delta^c
 =\tfrac i2 f_{abc}x_i^a\gamma^i_{\beta\delta}
                   \theta_\beta^b\theta_\delta^c,
\]
where the last equality uses antisymmetry of $f$ and relabels the
indices. Integration by parts now proves \eqref{eq:form} on $\Core$.

Whenever $f_{abc}\ne0$, the generators
$\theta_\beta^b$ and $\theta_\delta^c$ anticommute. Their product is
anti-Hermitian, so the displayed $B$ is Hermitian. It is a linear
function of $x$ with coefficients in the finite-dimensional space
$\operatorname{End}(\F)$. This proves \eqref{eq:fermion-bound}.
\end{proof}

Let $G=\Spin(9)$, and write $R_g\in\SO(9)$ for the vector
representation. Clifford multiplication gives a real orthogonal spin
representation $S_g$ on $\R^{16}$ satisfying
\begin{equation}
 S_g\gamma(v)S_g^{-1}=\gamma(R_gv),
 \qquad \gamma(v)=v_i\gamma^i.
 \label{eq:gamma-covariance}
\end{equation}
For $w\in\R^{16}$, write
$\theta^a(w)=w_\alpha\theta_\alpha^a$.

\begin{lemma}
\label{lem:rotations}
There is a unitary representation $U$ of $G$ on $\F$ such that
\[
 U_g\theta^a(w)U_g^{-1}=\theta^a(S_gw).
\]
It commutes with the gauge action. The representation
\begin{equation}
 (D_g\Psi)(x)=U_g\Psi(R_g^{-1}x)
 \label{eq:rotation-action}
\end{equation}
on $\Hphys$ preserves $\Core$ and the closed form $q$.

For an irreducible $G$ representation of type $\lambda$, let
$\Hphys_\lambda$ be its isotypic subspace and $P_\lambda$ the
orthogonal projection onto that subspace. Then $P_\lambda\Core$ is a
form core for $q$ on $\Hphys_\lambda$, and $\Hphys_\lambda$ reduces
$H$. The operator associated with the restricted form is the part
$H_\lambda=H|_{\Dom(H)\cap\Hphys_\lambda}$.
\end{lemma}

\begin{proof}
The action $S_g$ on each of the three color copies gives a homomorphism
$G\to\SO(48)$. Since $G$ is simply connected, it lifts to
$\Spin(48)$. Acting on $\F$ by the normalized Clifford generators
$\sqrt2\,\theta_\alpha^a$ gives $U$ with the stated property.

The actions of $G$ and $\SU(2)$ on the generator labels commute.
Their unitary commutator on $\F$ therefore commutes with every Clifford
generator and is scalar by irreducibility. For fixed $g\in G$ this
scalar, as a function of the gauge transformation, is a continuous
character of $\SU(2)$. Such a character is trivial. Thus the actions
commute on $\F$, and \eqref{eq:rotation-action} preserves $\Hphys$.

The pullback in \eqref{eq:rotation-action} sends the coordinate and
derivative vectors to their $R_g^{-1}$ transforms. The column of
Clifford generators transforms by $S_g^{-1}$. Applying
\eqref{eq:gamma-covariance}, and its consequence for $\gamma^{ij}$,
to the two terms in \eqref{eq:supercharges} gives
\[
 D_g Q_\alpha D_g^{-1}
   =(S_g^{-1})_{\alpha\beta}Q_\beta
 \quad\hbox{on }\Core.
\]
The orthogonality of $S_g$ proves $q(D_g\Psi)=q(\Psi)$.
Rotations preserve smoothness, and the rotated supports of a compact
set lie in a common ball, so $D_g\Core=\Core$.

The action is strongly continuous for the form norm
$\norm{\Psi}_q^2=\norm{\Psi}_2^2+q(\Psi)$. On $\Core$ this follows
directly from smoothness, \eqref{eq:supercharges}, and compact support;
it extends to the closure because every $D_g$ is a form-norm isometry.
For normalized Haar measure and the irreducible character $\chi_\lambda$,
\[
 P_\lambda
   =(\dim V_\lambda)\int_G\overline{\chi_\lambda(g)}D_g\,dg.
\]
This integral is bounded in form norm. Differentiation under the
integral and the common support bound show that it maps $\Core$ into
$\Core$. If $\Psi\in\Dom(q)\cap\Hphys_\lambda$, project any core
approximation to $\Psi$ to obtain a form-norm approximation in
$P_\lambda\Core$. This proves the core assertion.

Invariance of the closed form implies that its resolvent commutes with
every $D_g$: this follows, for example, from uniqueness of the solution
of
\[
 q(u,v)+\ip{u}{v}=\ip{f}{v}\qquad(v\in\Dom(q))
\]
defining $u=(H+1)^{-1}f$, where $q(u,v)$ denotes the polarization of
$q$. The resolvent therefore commutes with
$P_\lambda$. This proves reduction and identifies the restricted
operator with the stated part of $H$.
\end{proof}

\section{Transverse confinement in large rotation types}
\label{sec:confinement}

Fix a color $a$, a vector $u\in\R^9\setminus\{0\}$, and two real
numbers $s_b$ for $b\ne a$. Put $\widehat u=u/|u|$. For
$z=(z_b)_{b\ne a}\in(u^\perp)^2$, define $x(u,s,z)$ by its color
components:
\begin{equation}
 x(u,s,z)^a=u,\qquad
 x(u,s,z)^b=s_b\widehat u+z_b\quad(b\ne a).
 \label{eq:slice}
\end{equation}
The transverse space $(u^\perp)^2$ has dimension sixteen. On this
slice, the potential satisfies
\begin{equation}
 V(x(u,s,z))
 \geq\sum_{b\ne a}|u\wedge x(u,s,z)^b|^2
 =|u|^2|z|^2,\qquad |z|^2=\sum_{b\ne a}|z_b|^2.
 \label{eq:slice-potential}
\end{equation}
The subgroup $K_u\subset G=\Spin(9)$ fixing $u$ acts on the
transverse pair and is conjugate to $\Spin(8)$. We will compare the
$K_u$ types on a slice with those in each transverse oscillator level.
The needed lower bound on slice weights comes from restricting a
$G$ type to this stabilizer.

First use the reference stabilizer
$K=\operatorname{Stab}_G(e_9)\simeq\Spin(8)$.
Choose weight coordinates $\epsilon_1,\ldots,\epsilon_4$ for rotations
in four orthogonal planes in $\R^8=e_9^\perp$. The positive roots of
$K$ are $\epsilon_j\pm\epsilon_l$ for $j<l$; those of $G$ include
these and the $\epsilon_j$. Dominant weights have the forms
\[
 \lambda_1\geq\lambda_2\geq\lambda_3\geq\lambda_4\geq0
 \quad(G),\qquad
 \nu_1\geq\nu_2\geq\nu_3\geq|\nu_4|
 \quad(K),
\]
with coordinates all integral or all half-integral in each weight.

\begin{lemma}[Orthogonal restriction]
\label{lem:branching}
If the $K$ type of highest weight $\nu$ occurs in the restriction of
the $G$ type of highest weight $\lambda$, then
\begin{equation}
 \nu_1\geq\lambda_2.
 \label{eq:branching}
\end{equation}
\end{lemma}

\begin{proof}
This is the part of the classical orthogonal interlacing rule needed
here; compare \cite[Proposition~3.7]{BisiZygouras}. We give a character
proof that includes the half-integral cases.

Use formal torus monomials $Z^w$, allowing half-integral exponents, and
write $\Delta_G,\Delta_K$ for the Weyl denominators. The half-sums of
positive roots are
\[
 \rho_K=(3,2,1,0),\qquad
 \rho_G=\rho_K+(\tfrac12,\tfrac12,\tfrac12,\tfrac12).
\]
Thus
$\Delta_G=\Delta_K\prod_i(Z_i^{1/2}-Z_i^{-1/2})$.
Set $L_j=\lambda_j+4-j$. The Weyl character formula for type $B_4$
gives
\begin{equation}
 \chi_\lambda\Delta_K
 =\det\left(
   \frac{Z_i^{L_j+1/2}-Z_i^{-L_j-1/2}}
        {Z_i^{1/2}-Z_i^{-1/2}}
       \right)_{i,j=1}^4
 =\det\left(\sum_{s=-L_j}^{L_j}Z_i^s\right)_{i,j=1}^4,
 \label{eq:branch-character}
\end{equation}
where $s$ increases in steps of one.

Let $m_\eta$ be the multiplicity of the $K$ type $\eta$ in the
restriction, and let $W_K$ be the Weyl group of $K$. The Weyl
character formula for $K$ gives the second expression
\[
 \chi_\lambda\Delta_K
 =\sum_\eta m_\eta\sum_{w\in W_K}\det(w)\,
       Z^{w(\eta+\rho_K)}.
\]
Put $t=\nu+\rho_K$. Its coordinates satisfy
$t_1>t_2>t_3>|t_4|$, so $t$ is regular dominant for type $D_4$,
including when $t_4=0$. The coefficient of $Z^t$ in this expression
is therefore $m_\nu$: only $\eta=\nu$ and the identity Weyl element
contribute. In the determinant in \eqref{eq:branch-character}, row
$i$ involves only $Z_i$, so coefficient extraction acts row by row.
If the coordinate lattices of $\lambda$ and $\nu$ differ, the
coefficient is zero. Otherwise it is
\[
 \det\bigl(\mathbf1_{\{|t_i|\leq L_j\}}\bigr)_{i,j=1}^4.
\]
The $L_j$ strictly decrease, and so do the $|t_i|$. Each row is an
initial string of ones followed by zeros, and the string lengths are
nondecreasing down the rows. A nonzero determinant has neither a zero
row nor two equal rows, so these four lengths must be $1,2,3,4$.
In particular,
$t_1>L_2$, or $\nu_1>\lambda_2-1$. The difference
$\nu_1-\lambda_2$ is integral because the lattices agree. Hence
$\nu_1\geq\lambda_2$.
\end{proof}

For each $u$, choose a rotation carrying $e_9$ to $\widehat u$.
It conjugates $K$ to $K_u$; transport the torus and positive roots
by this conjugation. Lemma~\ref{lem:branching} then holds for $K_u$
in the transported weight coordinates.

\begin{lemma}[Types on a transverse slice]
\label{lem:slice}
Let $\Psi\in\Core\cap\Hphys_\lambda$. For the fixed parameters in
\eqref{eq:slice}, the function
\[
 \Phi(z)=\Psi(x(u,s,z))
\]
belongs, as a vector in $L^2((u^\perp)^2;\F)$, only to $K_u$
types whose highest weights satisfy $\nu_1\geq\lambda_2$.
The $K_u$ action here rotates $z$ and acts by $U$ on $\F$.
\end{lemma}

\begin{proof}
To restrict at fixed parameters, we first use the finite-dimensional
span of the smooth translates of $\Psi$.
The isotypic representation has the form
$\Hphys_\lambda\simeq V_\lambda\widehat\otimes\mathcal M_\lambda$,
where $\mathcal M_\lambda$ is a multiplicity Hilbert space. Writing
$d=\dim V_\lambda$, any individual vector has an expansion
$\sum_{j=1}^d e_j\otimes m_j$. Its span under $G$ is contained in
$V_\lambda\otimes\operatorname{span}\{m_1,\ldots,m_d\}$ and is
therefore finite dimensional. Apply this to $\Psi$, and denote its
translate span by $W_\Psi$.

Every vector in $W_\Psi$ has a smooth compactly supported
representative, because it is a finite linear combination of translates
of $\Psi$. An equality between these representatives in $L^2$ holds
pointwise by continuity. Thus restriction at the fixed $u,s_b$ is a
well-defined linear map from $W_\Psi$ to
$C_c^\infty((u^\perp)^2;\F)$. It intertwines $K_u$, which fixes
$u$ and the longitudinal coordinates $s_b$. Its image can have only
types occurring in $W_\Psi|_{K_u}$. Lemma~\ref{lem:branching} gives
the claimed bound. The $K_u$ action on the fiber $L^2$ space is
unitary, so distinct isotypic subspaces there are orthogonal.
\end{proof}

Choose $M\geq0$ so that the first coordinate of every $K$ weight of
$\F$ is at most $M$. Such an $M$ exists because $\F$ is finite
dimensional. The same bound holds for $K_u$ in the transported
coordinates.

The factor $1/4$ in the next estimate uses half of the kinetic term
in \eqref{eq:form}; the remaining quarter will provide the local
$H^1$ bound needed for compactness.

\begin{lemma}[Transverse oscillator bound]
\label{lem:oscillator}
Let $m\geq0$ be an integer and suppose
\begin{equation}
 \lambda_2>m+M.
 \label{eq:large-weight}
\end{equation}
For every slice $\Phi$ in Lemma~\ref{lem:slice},
\begin{equation}
 \int_{(u^\perp)^2}
   \left(\tfrac14|\nabla_z\Phi|^2+|u|^2|z|^2\norm{\Phi}^2\right)\,dz
 \ \geq\ (m+9)|u|\int_{(u^\perp)^2}\norm{\Phi}^2\,dz.
 \label{eq:fiber-oscillator}
\end{equation}
\end{lemma}

\begin{proof}
For $r=|u|>0$, the oscillator
\[
 h_r=-\tfrac14\Delta_z+r^2|z|^2
\]
has energies $r(n+8)$, $n=0,1,\ldots$. Indeed, the change of variables
$\xi=\sqrt{2r}\,z$ gives
$h_r=\frac r2(-\Delta_\xi+|\xi|^2)$ in sixteen dimensions.
Each vector at level $n$ is an invariant Gaussian times a polynomial
of total degree at most $n$. Tensor this level with $\F$.

The transverse coordinates are two copies of the vector representation
of $K_u$, whose weights are $\pm\epsilon_1,\ldots,\pm\epsilon_4$.
Every weight in a polynomial of degree at most $n$ has first
coordinate at most $n$, and every weight in the level tensored with
$\F$ consequently has first coordinate at most $n+M$. In particular,
the first coordinate of the highest weight of any constituent is at
most $n+M$.

Let $\Pi_n$ be the projection onto level $n$ tensored with $\F$.
It commutes with $K_u$ and hence with every $K_u$-isotypic projection
on the fiber $L^2$ space. By Lemma~\ref{lem:slice}, $\Pi_n\Phi$ can
therefore contain only types with $\nu_1\geq\lambda_2$. Since it
lies in level $n$ tensored with $\F$, its types also satisfy
$\nu_1\leq n+M$. For $n\leq m$ these bounds are incompatible by
\eqref{eq:large-weight}, so $\Pi_n\Phi=0$. Expanding in the complete
Hermite basis, the least remaining energy is $r(m+9)$, which proves
\eqref{eq:fiber-oscillator}.
\end{proof}

\begin{proposition}[Confinement in a rotation sector]
\label{prop:confinement}
Let $C$ be as in \eqref{eq:fermion-bound}, and choose an integer
$m\geq0$ such that
$c_m=(m+9)/3-C>0$. If $\lambda_2>m+M$, every
$\Psi\in\Dom(q)\cap\Hphys_\lambda$ belongs to
$H^1(\R^{27};\F)$, satisfies $|x|^{1/2}\Psi\in L^2(\R^{27};\F)$,
and obeys
\begin{equation}
 q(\Psi)\ \geq\ \tfrac14\norm{\nabla\Psi}_2^2
  +c_m\int_{\R^{27}}|x|\norm{\Psi(x)}^2\,dx.
 \label{eq:confinement}
\end{equation}
The operator $H_\lambda$ has compact resolvent and
$\ker H_\lambda=\{0\}$.
\end{proposition}

\begin{proof}
First take $\Psi\in\Core\cap\Hphys_\lambda$ and fix a color $a$.
For fixed $u$, the decomposition of the two other color variables
into $s_b,z_b$ is an orthogonal linear change of variables with unit
Jacobian. The transverse gradient uses only derivatives in those two
variables, so $|\nabla_z\Phi|^2\leq|\nabla\Psi|^2$. No derivative of
the decomposition with respect to $u$ is taken. Fubini's theorem,
using local orthogonal frames for $u^\perp$, therefore applies to
\eqref{eq:fiber-oscillator}; the set $u=0$ has measure zero.
Together with \eqref{eq:slice-potential}, it gives
\begin{equation}
 \tfrac14\norm{\nabla\Psi}_2^2+
   \int V(x)\norm{\Psi(x)}^2\,dx
 \ \geq\ (m+9)\int |x^a|\norm{\Psi(x)}^2\,dx.
 \label{eq:one-color}
\end{equation}
Average \eqref{eq:one-color} over the three colors; its left side
is unchanged. Splitting the kinetic term in \eqref{eq:form} into
two quarters and using \eqref{eq:fermion-bound} then gives
\[
 q(\Psi)\geq\tfrac14\norm{\nabla\Psi}_2^2+
 \int_{\R^{27}}\left(\tfrac{m+9}{3}\sum_a|x^a|-C|x|\right)
       \norm{\Psi(x)}^2\,dx.
\]
Since $\sum_a|x^a|\geq|x|$, this proves \eqref{eq:confinement}
on the core.

We now pass to the restricted form domain. By
Lemma~\ref{lem:rotations}, any vector in that domain is a form-norm
limit of vectors in
$\Core\cap\Hphys_\lambda$. Applied to differences of these vectors,
\eqref{eq:confinement} shows convergence of their gradients in $L^2$
and of their products with $|x|^{1/2}$ in $L^2$. The gradient and
multiplication by $|x|^{1/2}$ are closed operators. Their limits are
therefore the gradient and weighted product of the form limit.
Passing to the limit proves \eqref{eq:confinement} on the full
restricted form domain.

Its form-norm unit ball is bounded in $H^1(\R^{27};\F)$. Moreover,
for $R>0$, \eqref{eq:confinement} gives the uniform tail estimate
\[
 \int_{|x|>R}\norm{\Psi(x)}^2\,dx
 \leq\frac{q(\Psi)}{c_m R}.
\]
Rellich compactness on bounded balls applies componentwise in the
finite-dimensional space $\F$. A diagonal subsequence and the tail
estimate give compactness in $L^2(\R^{27};\F)$. Hence the inclusion
of the restricted form domain into $\Hphys_\lambda$ is compact.
The resolvent $(H_\lambda+1)^{-1}$ maps the Hilbert space boundedly
into that form domain, so it is compact.

If $\Psi\in\ker H_\lambda$, then $q(\Psi)=0$. The positive weighted
term in \eqref{eq:confinement} vanishes, forcing $\Psi=0$ because
$|x|>0$ almost everywhere. This proves the kernel assertion.
\end{proof}

\section{Physical rotation sectors and positive eigenvalues}
\label{sec:sectors}

It remains to produce an infinite-dimensional physical sector meeting
\eqref{eq:large-weight}. We first obtain a gauge-invariant fermionic
vector, then multiply a rotation highest vector by gauge-invariant
spatial polynomials.

\begin{lemma}
\label{lem:fermion-invariant}
The subspace $\F^{\SU(2)}$ is nonzero and invariant under $G$.
\end{lemma}

\begin{proof}
For $1\leq j\leq8$ and $1\leq a\leq3$, put
\[
 c_j^a=\frac{\theta_{2j-1}^a+i\theta_{2j}^a}{\sqrt2}.
\]
These twenty-four operators and their adjoints satisfy the canonical
anticommutation relations
$\{c_j^a,c_k^b\}=0$ and
$\{c_j^a,(c_k^b)^*\}=\delta_{jk}\delta_{ab}$.
The occupation operators $N_j^a=(c_j^a)^*c_j^a$ are commuting
projections. Start with a nonzero joint eigenvector. Applying an
occupied mode's annihilator removes that occupation and preserves all
others. The result is nonzero, since
$\norm{c_j^a v}^2=\ip{v}{N_j^a v}=\norm{v}^2$. Repeating this gives
a nonzero vector
annihilated by every $c_j^a$. For each subset of modes, apply its
creation operators in a fixed order. The resulting vectors are
nonzero and have distinct occupation patterns, hence are orthogonal.
Their span is invariant under every creation and
annihilation operator, so irreducibility makes it all of $\F$.
The common kernel of the annihilators is therefore exactly the
original line.

Gauge transformations mix the color indices and preserve the span
of the annihilators for each $j$. They therefore preserve the common
kernel line. Their action on that line is a continuous character of
$\SU(2)$ and is trivial. Thus $\F^{\SU(2)}\ne0$. Its $G$
invariance follows from the commutation in Lemma~\ref{lem:rotations}.
\end{proof}

\begin{lemma}
\label{lem:large-sectors}
There are $G$ types $\lambda$ with arbitrarily large $\lambda_2$
for which $\Core\cap\Hphys_\lambda$ is infinite dimensional.
\end{lemma}

\begin{proof}
Choose a nonzero highest vector $v$ of some irreducible $G$ type
$\mu$ in the finite-dimensional representation $\F^{\SU(2)}$.

We construct a gauge-invariant polynomial that raises the first two
highest-weight coordinates. In the complexified vector representation, set
\[
 w_1=e_1+i e_2,\qquad w_2=e_3+i e_4,
\]
choosing the torus convention in which their weights are
$\epsilon_1,\epsilon_2$. Use the complex bilinear extension of the
Euclidean dot product, and put
\[
 \zeta_a=w_1\cdot x^a,\qquad \eta_a=w_2\cdot x^a,\qquad
 P(x)=\sum_{a<b}(\zeta_a\eta_b-\zeta_b\eta_a)^2.
\]
The identity
\[
 P=(\zeta\cdot\zeta)(\eta\cdot\eta)-(\zeta\cdot\eta)^2
\]
uses bilinear dot products on the color space $\C^3$. The real
orthogonal adjoint action of $\SU(2)$ preserves these products,
so $P$ is gauge invariant. At $x=t(e_1,e_3,0)$ with $t>0$, one has
$P(x)=t^4$ and $|x|=\sqrt2\,t$. Thus $P$ is nonzero on every
nonempty radial shell.

For the action on spatial polynomials by pullback,
\[
 (w\cdot R_g^{-1}x)=(R_gw)\cdot x.
\]
Thus $w\mapsto(w\cdot x^a)$ intertwines the complexified vector
representation with linear polynomials, with the same weight sign.
The analogous map on exterior squares sends $w_1\wedge w_2$ to
$\zeta_a\eta_b-\zeta_b\eta_a$. The bivector $w_1\wedge w_2$ is
highest of weight $\epsilon_1+\epsilon_2$: a positive-root operator
either kills both factors or sends $w_2$ to a multiple of $w_1$,
and in the latter case its action on the wedge also vanishes.
Hence each minor is highest of that weight. Root operators act as
derivations on polynomials, so $P$ is highest of weight
$2(\epsilon_1+\epsilon_2)$. Its nonzero power $P^k$ is highest of
weight $2k(\epsilon_1+\epsilon_2)$.

Let $\chi\in C_c^\infty((0,\infty))$ be nonzero. The function
\[
 \Psi_\chi(x)=\chi(|x|)P(x)^k v
\]
is in $\Core$: it is smooth because its support avoids the origin,
and its factors are gauge invariant. The preceding ray shows that
$\Psi_\chi\ne0$. Multiplication by the radial function is $G$
equivariant. The translates of $\Psi_\chi$ lie in
the finite-dimensional representation obtained from homogeneous
polynomials of degree $4k$ tensored with $\F$. The vector is
highest of the dominant weight
\[
 \lambda=2k(\epsilon_1+\epsilon_2)+\mu
       =(2k+\mu_1,\,2k+\mu_2,\,\mu_3,\,\mu_4).
\]
In a decomposition of that finite representation into irreducibles,
each nonzero component of this vector is a weight-$\lambda$ vector
killed by all positive-root operators. It must be a highest vector,
so its summand has type $\lambda$. Therefore
$\Psi_\chi\in\Hphys_\lambda$.

Choosing $\chi$ with pairwise disjoint radial supports gives infinitely
many linearly independent vectors in the same sector. Finally,
$\lambda_2=2k+\mu_2$ tends to infinity with $k$.
\end{proof}

\begin{proof}[Proof of Theorem~\ref{thm:main}]
Choose an integer $m\geq0$ with $c_m=(m+9)/3-C>0$.
Lemma~\ref{lem:large-sectors} supplies a type $\lambda$ such that
$\lambda_2>m+M$ and $\Hphys_\lambda$ is infinite dimensional.
By Proposition~\ref{prop:confinement}, $H_\lambda$ has compact
resolvent and no zero eigenvector.

Since $H_\lambda$ is nonnegative with compact resolvent on an
infinite-dimensional Hilbert space, the compact-resolvent spectral
theorem gives an orthonormal eigenbasis $(\Psi_j)$ with eigenvalues
$E_j\to\infty$. Its trivial kernel makes every $E_j$ positive.
By Lemma~\ref{lem:rotations}, $H_\lambda$ is the part of $H$ on a
reducing subspace. These vectors therefore belong to $\Dom(H)$ and
satisfy $H\Psi_j=E_j\Psi_j$ in the full physical Hilbert space.
\end{proof}

\begin{thebibliography}{99}

\bibitem{ArefevaKoshelevMedvedev}
I. Ya. Aref'eva, A. S. Koshelev, and P. B. Medvedev,
\emph{On stable sector in supermembrane matrix model},
Nuclear Physics B \textbf{579} (2000), 411--436.
\url{https://doi.org/10.1016/S0550-3213(00)00205-4}.

\bibitem{BFSS}
Tom Banks, Willy Fischler, Stephen H. Shenker, and Leonard Susskind,
\emph{M theory as a matrix model: A conjecture},
Physical Review D \textbf{55} (1997), 5112--5128.
\url{https://doi.org/10.1103/PhysRevD.55.5112}.

\bibitem{BisiZygouras}
Elia Bisi and Nikos Zygouras,
\emph{Transition between characters of classical groups, decomposition of
Gelfand--Tsetlin patterns and last passage percolation},
arXiv:1905.09756, version 5, 2022.
\url{https://arxiv.org/abs/1905.09756}.

\bibitem{DanielssonFerrettiSundborg}
Ulf H. Danielsson, Gabriele Ferretti, and Bo Sundborg,
\emph{D-particle dynamics and bound states},
International Journal of Modern Physics A \textbf{11} (1996), 5463--5478.
\url{https://doi.org/10.1142/S0217751X96002492}.

\bibitem{DaviesSimon}
E. B. Davies and Barry Simon,
\emph{Spectral properties of Neumann Laplacian of horns},
Geometric and Functional Analysis \textbf{2} (1992), 105--117.
\url{https://doi.org/10.1007/BF01895707}.

\bibitem{deWitLuscherNicolai}
B. de Wit, M. L\"uscher, and H. Nicolai,
\emph{The supermembrane is unstable},
Nuclear Physics B \textbf{320} (1989), 135--159.
\url{https://doi.org/10.1016/0550-3213(89)90214-9}.

\bibitem{Luscher}
Martin L\"uscher,
\emph{Some analytic results concerning the mass spectrum of Yang--Mills
gauge theories on a torus},
Nuclear Physics B \textbf{219} (1983), 233--261.
\url{https://doi.org/10.1016/0550-3213(83)90436-4}.

\bibitem{SethiSternInvariance}
Savdeep Sethi and Mark Stern,
\emph{Invariance theorems for supersymmetric Yang--Mills theories},
Advances in Theoretical and Mathematical Physics \textbf{4} (2000), 487--501.
\url{https://doi.org/10.4310/ATMP.2000.v4.n2.a8}.

\bibitem{Simon}
Barry Simon,
\emph{Some quantum operators with discrete spectrum but classically
continuous spectrum},
Annals of Physics \textbf{146} (1983), 209--220.
\url{https://doi.org/10.1016/0003-4916(83)90057-X}.

\bibitem{Smilga}
A. V. Smilga,
\emph{Comments on thermodynamics of supersymmetric matrix models},
Nuclear Physics B \textbf{818} (2009), 101--114.
\url{https://doi.org/10.1016/j.nuclphysb.2009.03.023}.

\end{thebibliography}

\end{document}
